\documentclass[11pt,a4paper]{article}

\usepackage[utf8]{inputenc}
\usepackage[T1]{fontenc}
\usepackage[margin=2.3cm]{geometry}
\usepackage{amsmath,amssymb}
\usepackage{listings}
\usepackage{verbatim}
\usepackage{xcolor}
\usepackage{hyperref}
\usepackage{booktabs}

\hypersetup{colorlinks=true, linkcolor=blue, urlcolor=blue, citecolor=blue}

\definecolor{lstbg}{rgb}{0.97,0.97,0.97}
\definecolor{lstframe}{rgb}{0.70,0.70,0.70}

\lstset{
    basicstyle=\ttfamily\footnotesize,
    keywordstyle=\color{blue}\bfseries,
    commentstyle=\color{gray}\itshape,
    stringstyle=\color{orange},
    showstringspaces=false,
    breaklines=true,
    frame=single,
    rulecolor=\color{lstframe},
    backgroundcolor=\color{lstbg},
    tabsize=4,
    captionpos=b,
}

\lstdefinestyle{plain}{
    basicstyle=\ttfamily\footnotesize,
    breaklines=true,
    frame=single,
    rulecolor=\color{lstframe},
    backgroundcolor=\color{lstbg},
    keywordstyle={},
    commentstyle={},
    stringstyle={},
}

\title{\textbf{Report delle modifiche a \texttt{cssp.py}}\\[0.3em]
\large Simulazione MNRS Quantum Walk Search per il CSSP:\\
cronologia degli interventi e diagnosi finale}
\author{Jacopo Bellosi}
\date{Aprile 2026}

\begin{document}
\maketitle

% =====================================================================
\section{Il problema di partenza}
% =====================================================================

Il paper \emph{``Solving the Subset Sum Problem via Quantum Walk
Search''} (Lancellotti et al., IEEE TC, vol.~75, n.~1, 2026) fornisce
un'implementazione di riferimento in \texttt{myQLM} nel repository
GitHub \texttt{polimi-dl-lab/paper-codes/2025-TC}. L'assunto iniziale
era che tale codice potesse essere eseguito end-to-end su un simulatore
classico (\texttt{PyLinalg}) producendo un'amplificazione osservabile
della soluzione nel registro dei vertici $S_1$. L'esecuzione diretta di
\texttt{cssp.py} sull'istanza $n{=}3$, $k{=}1$,
\texttt{values}=[1,2,3], \texttt{target\_sum}=3 ha invece prodotto una
distribuzione di probabilità \emph{piatta} sui tre valori ammissibili.
Un secondo run con checkpoint intermedi mostra che il circuito viene
costruito e simulato correttamente, ma la fase marcata dall'oracolo non
viene trasformata in ampiezza dal passo di diffusione.

\begin{table}[h]
\centering
\small
\begin{tabular}{@{}p{0.47\linewidth} p{0.47\linewidth}@{}}
\toprule
\textbf{(a) Esecuzione diretta di \texttt{cssp.py}} &
\textbf{(b) Esecuzione con checkpoint intermedi} \\
\midrule
\begin{minipage}[t]{\linewidth}
\begin{lstlisting}[style=plain]
n=3, k=1, values=[1,2,3], target=3
{'nbqbits': 28, 'gate_size': 2547}

0.33333333333 |1>
0.33333333333 |2>
0.33333333333 |3>
\end{lstlisting}
\end{minipage}
&
\begin{minipage}[t]{\linewidth}
\begin{lstlisting}[style=plain]
Checkpoint 1 (Dicke + VBE)   : uniform
Checkpoint 2 (Edge superp.)  : 12 stati, p=1/12
Checkpoint 3 (Post-Oracle)   : phase flip OK
Checkpoint 4 (Post-QPE/QFT)  : stati sparsi
Checkpoint 5 (Post-Diff.)    : nessuna amp.
Checkpoint 6 (Fine iterazione): stato misto

0.33333333333 |1>
0.33333333333 |2>
0.33333333333 |3>
\end{lstlisting}
\end{minipage}\\
\bottomrule
\end{tabular}
\caption{Output dell'esecuzione originale: la distribuzione finale su
$S_1$ coincide con la distribuzione iniziale del Dicke state
$|D^3_1\rangle$. L'oracolo applica correttamente il phase flip agli
stati marcati (Checkpoint~3), ma la riflessione approssimata non
converte la fase in un'ampiezza osservabile.}
\label{tab:baseline}
\end{table}

In nessuno dei run osservati il simulatore produce errori: il circuito
viene compilato, linkato e simulato fino in fondo. Il problema non è
dunque di natura implementativa ``di basso livello'' ma algoritmica o
di regime: le modifiche descritte nella Sezione~\ref{sec:mod} sono
stati i tentativi di fare emergere l'amplificazione attesa.

% =====================================================================
\section{Le modifiche (ordine cronologico da \texttt{git log})}
\label{sec:mod}
% =====================================================================

La cronologia che segue è ricostruita dallo storico dei commit del
repository locale; commit multipli con la stessa idea di fondo sono
stati raggruppati in un'unica sottosezione. Tutte le misure finali
sono state prese sulla stessa istanza canonica $J(3,1)$ con target
$p{=}3$, così da tenere una baseline omogenea.

% ---------------------------------------------------------------------
\subsection*{Modifica 1 -- Correzione della logica di Reflection B
(commit \texttt{475436d}, 17 marzo 2026)}
% ---------------------------------------------------------------------

Nel codice originale la seconda riflessione del walk operatore,
$U_{\text{ref}}(a_{v_j})$, era implementata in forma asimmetrica
rispetto alla prima: il phase flip era controllato solo sul registro
$\omega'$ (\texttt{wstate\_zeros}) e gli $X$--gate di specchiamento
erano applicati solo parzialmente. La Ref~B deve invece riflettere
attorno allo stato \emph{coin} completo $|0\rangle^{\otimes n}$, quindi
serve un $Z$ controllato sull'intero registro coin $\omega\cup\omega'$
e gli $X$--gate vanno applicati simmetricamente su entrambi i
sottoregistri.

\begin{lstlisting}[language=Python]
# --- Prima (originale) ---
prw.apply(qrw_update.dag(),
          node_s_zeros, node_s_ones,
          node_t_zeros, node_t_ones,
          alpha_zeros, alpha_ones,
          wstate_zeros, wstate_ones)
for j in range(n - k):
    prw.apply(X, wstate_zeros[j])
prw.apply(Z.ctrl(n - k), qpe_s[qw_iter], wstate_zeros)
for j in range(k):   # BUG: dovrebbe essere range(n-k)
    prw.apply(X, wstate_zeros[j])

# --- Dopo (commit 475436d) ---
prw.apply(qrw_update.dag(),
          node_t_ones, node_t_zeros,
          node_s_ones, node_s_zeros,
          alpha_ones, alpha_zeros,
          wstate_ones, wstate_zeros)
for j in range(k):
    prw.apply(X, wstate_ones[j])
for j in range(n - k):
    prw.apply(X, wstate_zeros[j])
prw.apply(Z.ctrl(n), qpe_s[qw_iter],
          wstate_ones, wstate_zeros)    # controllo su intero coin
for j in range(k):
    prw.apply(X, wstate_ones[j])
for j in range(n - k):
    prw.apply(X, wstate_zeros[j])
prw.apply(qrw_update,
          node_t_ones, node_t_zeros,
          node_s_ones, node_s_zeros,
          alpha_ones, alpha_zeros,
          wstate_ones, wstate_zeros)
\end{lstlisting}

\begin{lstlisting}[style=plain]
n=3, k=1, values=[1,2,3], target=3
{'nbqbits': 28, 'gate_size': 2581}
0.33333333333 |1>
0.33333333333 |2>
0.33333333333 |3>
(runtime: 11h 37min 54s CPU)
\end{lstlisting}

La Reflection~B diventa ora strutturalmente corretta rispetto alla
descrizione del paper, ma l'output finale non cambia: distribuzione
ancora piatta. Il bug esisteva ma non era la causa radice.

% ---------------------------------------------------------------------
\subsection*{Modifica 2 -- Rimozione del blocco di cleanup
$\alpha/\omega$ (commit \texttt{7e4e8a3} + \texttt{44e3e0e},
20--23 marzo 2026)}
% ---------------------------------------------------------------------

All'interno del blocco \texttt{compute()} della QPE, dopo le $\ell_s$
applicazioni controllate del walk, il codice originale eseguiva una
``cleanup'' dei registri $\alpha$ e $\omega$: un ciclo di
\texttt{copy\_register(m).ctrl()} seguito da \texttt{generate(...)}. Il
ragionamento originario era che si trattasse di un \emph{uncompute}
degli ancilla per mantenere la reversibilità della routine. In realtà
quei registri al termine del ciclo non sono in $|0\rangle$ ma in
superposizione con i vertici visitati dal walk: invocare
\texttt{generate} come se fossero $|0\rangle$ aggiunge gate invece di
rimuoverli, collassando la coerenza del walk prima del QFT.

\begin{lstlisting}[language=Python]
# --- Prima (originale) ---
for j in range(k):
    prw.apply(
        qregs_init.copy_register(m).ctrl(),
        wstate_ones[j], node_s_ones[j], alpha_ones)
for j in range(n - k):
    prw.apply(
        qregs_init.copy_register(m).ctrl(),
        wstate_zeros[j], node_s_zeros[j], alpha_zeros)
prw.apply(generate(k, 1),       wstate_ones)
prw.apply(generate(n - k, 1),   wstate_zeros)

# --- Dopo (commit 7e4e8a3 + 44e3e0e) ---
# Blocco commentato/rimosso: distruggeva la coerenza del walk.
# L'uncompute verra' fatto implicitamente dal prw.uncompute()
# del blocco compute(), con i .dag() inversi delle gate che
# compongono la QPE.
\end{lstlisting}

\begin{lstlisting}[style=plain]
n=3, k=1, values=[1,2,3], target=3
{'nbqbits': 28, 'gate_size': 1359}     # -47% di gate
0.33333333333 |1>
0.33333333333 |2>
0.33333333333 |3>
\end{lstlisting}

Il numero di gate si riduce drasticamente, confermando che il blocco
era effettivamente superfluo/errato. Ma la distribuzione osservata non
cambia: il cleanup non era ciò che nascondeva l'amplificazione.

% ---------------------------------------------------------------------
\subsection*{Modifica 3 -- Inserimento \emph{low-width} al posto di
\emph{low-depth} (commit \texttt{21b2b0f}, 22 marzo 2026)}
% ---------------------------------------------------------------------

Nel data structure \emph{sliding-sort array} usato per
inserire/estrarre elementi da $S_1/S_0$, la variante
\texttt{insert\_ld} (\emph{low-depth}) lascia ancilla non azzerate
dopo il \texttt{.dag()}, rompendo la reversibilità del walk operatore.
La variante \texttt{insert\_lw} (\emph{low-width}) è reversibile, al
costo di un incremento modesto di depth. È stata introdotta sia per la
simulazione sia come default dell'operatore \texttt{delete}.

\begin{lstlisting}[language=Python]
# --- Prima ---
from .sliding_sort_array import insert_ld
insert = insert_ld                       # lascia garbage in ancilla

# --- Dopo (commit 21b2b0f) ---
from .sliding_sort_array import insert_lw, insert_ld
insert = insert_lw if low_width else insert_ld   # reversibile
# delete() aggiornato in sliding_sort_array.py per usare
#    insert_lw(...).dag() internamente
\end{lstlisting}

\begin{lstlisting}[style=plain]
n=3, k=1, values=[1,2,3], target=3
{'nbqbits': 28, 'gate_size': 1366}
0.33333333333 |1>
0.33333333333 |2>
0.33333333333 |3>
\end{lstlisting}

La routine di insert/delete diventa pulita rispetto agli ancilla (ora
il walk è reversibile come richiesto dalla teoria MNRS), ma questo
tocca la correttezza strutturale del $U_U$ senza modificare la
dinamica del sistema. Il risultato resta piatto.

% ---------------------------------------------------------------------
\subsection*{Modifica 4 -- Clamp dello spectral gap e soglia minima per
$\ell_s$ (commit \texttt{51799b8}, 27 marzo 2026)}
% ---------------------------------------------------------------------

La formula del paper per lo \emph{spectral gap} di $J(n,k)$ è
\[
\delta \;=\; \frac{n}{k(n-k)},
\qquad
\ell_s \;=\; \left\lceil \log_2\!\frac{\pi}{2\sqrt{\delta}} \right\rceil.
\]
Per l'istanza canonica $(n,k)=(3,1)$ si ottiene $\delta = 1.5$, valore
impossibile per uno spectral gap di una catena di Markov (che vive in
$(0,1]$), e $\ell_s = 1$, insufficiente per una QPE significativa. È
stato inserito un \emph{clamp} a $\delta \le 1$ e una soglia minima di
3 qubit per il registro QPE:

\begin{lstlisting}[language=Python]
# --- Prima ---
delta  = n / (k * (n - k))
len_s  = int(np.ceil(np.log2(np.pi / (2 * np.sqrt(delta)))))

# --- Dopo (commit 51799b8) ---
# delta in (0,1]: la formula del paper ha senso asintotico; per
# k(n-k) < n la espressione esplode e va sottoposta a ceiling.
delta  = min(n / (k * (n - k)), 1.0)
# un singolo qubit di QPE distingue solo fase 0 o pi: insufficiente
# per approssimare la riflessione nel subspazio stazionario.
len_s  = max(int(np.ceil(np.log2(np.pi / (2 * np.sqrt(delta))))), 3)
\end{lstlisting}

\begin{lstlisting}[style=plain]
n=3, k=1, values=[1,2,3], target=3
{'nbqbits': 28 -> 30, 'gate_size': 1366}
0.33333333333 |1>
0.33333333333 |2>
0.33333333333 |3>
\end{lstlisting}

La precisione della QPE aumenta (3 qubit anziché 1), ma la distribuzione
finale resta piatta. Motivo: il clamp è un patch numerico, non
introduce uno spectral gap che il grafo non possiede (vedi
Sezione~\ref{sec:diag}).

% ---------------------------------------------------------------------
\subsection*{Modifica 5 -- QFT inversa al posto della QFT diretta
(commit \texttt{bf13d21}, 13 aprile 2026)}
% ---------------------------------------------------------------------

Lo schema standard della Quantum Phase Estimation è
\[
\text{QPE} \;=\; (\mathrm{QFT}^{\dagger} \otimes I)\,
                 \prod_{j=0}^{\ell_s-1}
                 \Lambda_j\!\bigl(U_W^{2^{j}}\bigr)\,
                 (H^{\otimes \ell_s}\otimes I),
\]
dove $\mathrm{QFT}^{\dagger}$ è la QFT inversa. Nel codice del paper
veniva invece applicata la QFT diretta ($\mathrm{QFT}$), che nella
catena dei \texttt{compute()}/\texttt{uncompute()} di myQLM produce
uno schema non simmetrico rispetto alla definizione della QPE.

\begin{lstlisting}[language=Python]
# --- Prima ---
prw.apply(QFT(len_s), qpe_s)          # QFT diretta

# --- Dopo (commit bf13d21) ---
prw.apply(QFT(len_s).dag(), qpe_s)    # QFT inversa (come da QPE)
\end{lstlisting}

\begin{lstlisting}[style=plain]
n=3, k=1, values=[1,2,3], target=3
{'nbqbits': 28, 'gate_size': 2551}
0.33333333333 |1>
0.33333333333 |2>
0.33333333333 |3>
(runtime: 12h 30min 12s CPU)
\end{lstlisting}

Correzione teoricamente dovuta ma senza effetto misurabile sul
risultato: la QFT e la sua inversa coincidono nel caso $\ell_s=1$
(entrambe si riducono a $H$), e per $\ell_s=3$ la differenza si
manifesta in riflessi che il walk a questa scala non riesce comunque a
trasformare in ampiezza.

% ---------------------------------------------------------------------
\subsection*{Modifica 6 -- Debugger a checkpoint per l'ispezione delle
ampiezze (commit \texttt{1408452}, 17 aprile 2026)}
% ---------------------------------------------------------------------

Per capire dove esattamente si perdesse l'amplificazione è stata
creata una copia strumentata \texttt{cssp\_with\_checkpoint.py} che
simula lo stato globale in sei punti chiave (Dicke+VBE, edge
superposition, post-oracolo, post-QPE, post-diffusione, fine
iterazione). A ogni checkpoint si stampano probabilità e ampiezze
complesse di tutti gli stati con $p>5\cdot 10^{-4}$, opportunamente
suddivise per registro.

\begin{lstlisting}[language=Python]
def simulate_and_print_state(prw, n, k, m, len_s, n_qubits_sum, label):
    circ = prw.to_circ(link=[classarith, cuccaro_arith])
    res  = PyLinalg().submit(circ.to_job())
    for sample in res:
        if sample.probability > 5e-4:
            raw = str(sample.state)[1:-1]
            # ... parsing in dicke/S_1/S_0/t_*/a_*/w_*/qpe/sum ...
            print(f"Prob: {sample.probability:.5f} "
                  f"| Ampl: {sample.amplitude:+.4f}"
                  f"| S_1:{s1_str} | Sum:{sum_str}")
\end{lstlisting}

\begin{lstlisting}[style=plain]
Checkpoint 2 (Edge Superposition):
  Prob 0.0833 | +0.2887 | S_1:0 | Sum:2     (non-marcati)
  Prob 0.0833 | +0.2887 | S_1:1 | Sum:3     (marcati)   * 8 stati

Checkpoint 3 (Post-Oracle):
  Prob 0.0833 | +0.2887 | S_1:0 | Sum:2     (non-marcati)
  Prob 0.0833 | -0.2887 | S_1:1 | Sum:3     <== phase flip OK

Checkpoint 5 (Post-Diffusion):
  Prob 0.0417 ... ampiezze +/-0.2041 sparse
  (nessuna concentrazione su Sum:3)

Checkpoint 6 (End of Iteration):
  Prob 0.0833 ... ampiezze +/-0.2887 con segni misti su S_1=0 e S_1=1
\end{lstlisting}

Questa modifica non tocca l'algoritmo ma è stata decisiva per la
diagnosi: conferma empiricamente che \emph{l'oracolo funziona} (phase
flip sui marcati al Checkpoint~3) e che il problema nasce fra il
Checkpoint~3 e il Checkpoint~5, cioè all'interno della riflessione
approssimata QPE--based.

% =====================================================================
\section{Perché non funziona: la diagnosi finale}
\label{sec:diag}
% =====================================================================

\subsection{Il verdetto tecnico}

Dopo le sei modifiche della Sezione~\ref{sec:mod}, la distribuzione
finale su $S_1$ resta identicamente $(\tfrac13,\tfrac13,\tfrac13)$. La
causa non è una singola bug localizzata ma un disallineamento fra il
regime di validità dell'algoritmo MNRS e le istanze scelte come test.
Si distinguono tre casi.

\paragraph{Caso 1: $J(3,1)$ (l'istanza usata per il debug).}
Per $(n,k)=(3,1)$ la formula del paper restituisce
\[
\delta = \frac{3}{1\cdot 2} = \frac{3}{2} > 1,
\]
violazione dell'assunzione $\delta \in (0,1]$: lo spectral gap di una
catena di Markov reversibile è per costruzione al più $1$. La formula
$\delta = n/(k(n-k))$ è un'espressione \emph{asintotica} valida per
$k(n-k) \ge n$ (cioè $k \ge 2$ e $n \ge 4$), e degenera altrove.
Nessuna scelta di $\ell_s$, di direzione della QFT, di scaling
$U_W^{2^j}$ delle potenze del walk o di variante della Reflection~B
può produrre una distribuzione non--uniforme su $S_1$: la riflessione
approssimata $U_{\mathrm{ref}}(E)$ non riesce ad agire come $-I$ su
$|E\rangle$ con $\le 1$ qubit di QPE ``vero'' (prima del clamp
artificiale a $\ell_s=3$), e aumentare $\ell_s$ artificialmente non
può creare uno spectral gap che il grafo non ha.

\paragraph{Caso 2: $J(4,2)$.}
Per $(n,k)=(4,2)$ si ha $\delta = 4/(2\cdot 2) = 1$ esattamente, caso
di bordo. L'algoritmo è in principio definito ma senza separazione
spettrale da sfruttare: aumentare $\ell_s$ a 3--5 qubit e usare la
QFT inversa con scaling esponenziale $U_W^{2^j}$ produce una
amplificazione marginale, ma con solo
$N = \binom{4}{2} = 6$ vertici un Grover diretto sul registro dei
vertici è più semplice e strettamente migliore in termini di
probabilità di successo dopo un numero fissato di iterazioni.

\paragraph{Caso 3: $J(n,k)$ con $\delta<1$, $n\ge 5$, $k\ge 2$.}
Per istanze ``genuine'' (ad esempio $(5,2)$, $(6,2)$, $(6,3)$) il
circuito MNRS così come patchato nella Sezione~\ref{sec:mod} dovrebbe
amplificare la soluzione. Il vincolo pratico è di natura computazionale
(Sezione~\ref{subsec:cause}): il numero di qubit $\Theta(n\log n)$ e il
numero di gate del walk spingono il runtime di \texttt{PyLinalg} oltre
ciò che è tollerabile su un singolo nodo di calcolo. L'istanza $(3,1)$
era stata scelta proprio perché simulabile, non perché soddisfacesse
le precondizioni teoriche.

\subsection{Le cause distinte}
\label{subsec:cause}

È utile separare cause indipendenti che in superficie si confondono:

\begin{enumerate}
\item \textbf{Limiti fisici della simulazione classica.}
      Il circuito su $(3,1)$ usa già 28 qubit (Dicke+$S_1$+$S_0$+
      $T_1$+$T_0$+$\alpha_1$+$\alpha_0$+$\omega_1$+$\omega_0$+
      $\text{QPE}$+\text{sum}). \texttt{PyLinalg} è un simulatore a
      vettore di stato completo: ogni gate tocca un vettore di
      $2^{28} \approx 2.7\cdot 10^{8}$ ampiezze complesse. Con
      $\sim 2500$ gate per iterazione si arriva al limite di quello
      che un singolo core CPU può digerire in tempi dell'ordine delle
      ore. Aumentare $(n,k)$ per entrare nel regime $\delta<1$ porta a
      $\ge 35$ qubit e rende la simulazione impraticabile
      ($2^{35}\approx 3.4\cdot 10^{10}$ ampiezze).

\item \textbf{Limiti strutturali dell'algoritmo su istanze banali.}
      MNRS richiede $\delta<1$ con margine per concentrare la misura
      sullo stato marcato. $J(3,1)$ \emph{non soddisfa} questa
      precondizione: lo spectral gap reale è $\Theta(1)$ (il grafo è
      $K_3$, tre vertici tutti adiacenti), non c'è una direzione
      ``lenta'' lungo cui la QPE possa discriminare lo stato
      stazionario $|E\rangle$ dal resto. Tutte le tecniche
      diagnostiche a valle---Ref~B simmetrica, QFT inversa, scaling
      $2^j$, clamp di $\delta$---correggono bug reali ma non inventano
      uno spectral gap inesistente.

\item \textbf{Scelte di design del paper incompatibili con la
      simulazione classica.}
      Il codice \texttt{2025-TC} è stato progettato come banco di
      prova per misurare T-count e T-depth su circuiti di dimensione
      \emph{crittografica}; gli autori lavorano su una piattaforma
      Atos QLM (hardware emulato istituzionale) per tempi di
      \emph{counting} su grandi istanze, non per distribuzioni di
      probabilità su piccole istanze. La preparazione via Dicke +
      VBE, la doppia coppia di registri $S/T$ e la QPE esplicita sono
      tutte scelte ottimizzate per un regime dove le costanti
      moltiplicative sono irrilevanti, e che su istanze giocattolo
      generano overhead enormi a fronte di risultati qualitativi
      nulli.

\item \textbf{Bug veri nel codice (indipendenti dai limiti fisici).}
      La Modifica~1 (Ref~B asimmetrica) e la Modifica~2 (cleanup
      errato) sono bug reali: la prima violava la definizione di
      $U_W = U_{\mathrm{ref}}(a_{v_j})\,U_{\mathrm{ref}}(a_{v_i})$, la
      seconda rompeva la reversibilità del blocco
      \texttt{compute()/uncompute()}. Se il regime fosse stato corretto
      (grande $n$, $\delta<1$) questi bug avrebbero effettivamente
      soppresso l'amplificazione. Su $J(3,1)$, rimuoverli non ha
      effetto misurabile perché la macchina non era in un regime in
      cui qualcosa fosse da amplificare.
\end{enumerate}

\subsection{Le alternative percorribili}

\paragraph{Alternativa 1: Grover diretto sul registro dei vertici.}
Per $N = \binom{n}{k}$ piccolo ($\le 20$--30) si ignora l'apparato
MNRS e si costruisce un classico Grover:
$\lfloor\tfrac{\pi}{4}\sqrt{N}\rfloor$ iterazioni, un solo registro di
$\lceil\log_2 N\rceil$ qubit che indicizza le combinazioni possibili, e
un oracolo che valuta $\sum_{i\in S} x_i = p$. Funziona dove MNRS
fallisce perché non dipende dallo spectral gap del grafo: la
concentrazione di misura viene dalla struttura della riflessione
attorno alla media, non da una QPE che richiede un gap spettrale
osservabile. Il limite è che perde l'interesse teorico del paper
(nessuno vuole riprodurre il paper per scoprire che Grover basta su
istanze giocattolo).

\paragraph{Alternativa 2: MNRS semplificato con encoding compatto.}
Per un'istanza fissata come $(n,k) = (5,2)$ o $(6,3)$, si può:
\begin{itemize}
\item eliminare il registro Dicke e il doppio cover $S_0$ e i registri
      $T_1/T_0$ inizializzando direttamente $|S\rangle$ su un solo
      registro a $\lceil\log_2 N\rceil$ qubit;
\item rimpiazzare la QPE esplicita con una riflessione
      \emph{esatta} sul walk operatore (calcolabile in forma chiusa
      per grafi di Johnson piccoli);
\item usare $\omega$ come coin a un solo qubit per $k=1$ o mantenere
      il coin binario solo quando $k\ge 2$.
\end{itemize}
L'obiettivo non è più la simulazione fedele del paper ma la
dimostrazione qualitativa dell'amplificazione. Qubit count tipico:
$8\text{--}14$, raggiungibile da \texttt{PyLinalg} in minuti. Limite:
si rinuncia alla corrispondenza uno--a--uno con la costruzione del
paper, quindi il risultato non è una verifica del paper ma una
variante.

\paragraph{Alternativa 3: hardware QLM istituzionale.}
Il gruppo POLIMI che ha prodotto il paper lavora su un'istanza di Atos
QLM con simulatore MPS / \texttt{LinAlg} distribuito. In quel regime si
possono affrontare istanze $n\sim 10$, $k\in\{3,4\}$, dove
$\delta\lesssim 0.5$ e $\ell_s\ge 3$ è naturalmente rispettato, e il
runtime è di ore su un cluster anziché giorni su un laptop. Ottenere
accesso a quella macchina (o almeno a \texttt{qat-utils}, il pacchetto
interno con i backend distribuiti) renderebbe possibile una verifica
fedele del circuito del paper. Limite: è una risorsa istituzionale,
non riproducibile autonomamente; non sostituisce un fix del codice, ma
rende la simulazione fattibile fuori dall'istanza giocattolo dove
l'algoritmo semplicemente non può amplificare.

\medskip
\noindent\textbf{Conclusione.} La pipeline MNRS del paper è
algoritmicamente corretta (al netto dei due bug localizzati in
Modifica~1 e Modifica~2, ora corretti in \texttt{code\_updated/}) ma
non è eseguibile \emph{in forma significativa} sulle istanze
giocattolo $J(3,1)$ e $J(4,2)$ usate per il debug, perché tali istanze
cadono fuori dal regime $\delta<1$ in cui MNRS è progettato per
funzionare. Le alternative percorribili sono (a) rinunciare alla
fedeltà e usare Grover, (b) rinunciare alla simulazione fedele su
laptop e costruire una variante compatta, o (c) ottenere accesso
all'hardware su cui il paper è stato in realtà eseguito.

\end{document}
